\scanpage{12}\mspage{11}
\diagram{%
\[
\vcenter{\hbox{\begin{tikzcd}[ampersand replacement=\&,row sep=1.25em,column sep=.2em,
execute at end picture={
\draw ([xshift=-.45em,yshift=1em]leq.north west) -- ([xshift=.8em,yshift=.15em]leq.north east);
\draw ([xshift=.55em,yshift=.26em]leq.north east) -- ([xshift=.75em,yshift=-.95em]ltlt.south east) -- ([xshift=-.9em,yshift=-.1em]ltlt.south west);
\draw ([xshift=-.65em,yshift=.93em]leq.north west) -- ([xshift=-.84em,yshift=.08em]ltlt.south west);
}]
\& \& G\triangleleft F \arrow[d,Rightarrow] \& \& \\
G\preccurlyeq F \arrow[dr,Rightarrow] \& \& |[alias=leq]| G\leq F \arrow[dl,Rightarrow] \arrow[dr,Rightarrow] \& \& G\parallel F \arrow[dl,Rightarrow] \\
\& |[alias=ltlt]| G\ll F \& \& G\unsure{\mathrel{\tikz[baseline=0pt,x=1em,y=1em]{\draw[line width=.6pt] (-.38,.62)--(.38,.25)--(-.43,-.10); \draw[line width=.6pt] (.34,.62)--(-.32,.25)--(.38,-.10); \draw[line width=.6pt] (-.12,.67)--(-.12,-.08);}}|\mathrel{\tikz[baseline=0pt,x=1em,y=1em]{\draw[line width=.6pt] (-.38,.62)--(.38,.25)--(-.43,-.10); \draw[line width=.6pt] (.34,.62)--(-.32,.25)--(.38,-.10); \node[inner sep=0pt] at (-.12,.28) {$\deleted{\mid}$};}}} \unsure{F|\mathbb{F}} \&
\end{tikzcd}}} \tag*{(1.18)}
\]}

Among the six relations, there are four order relations $\triangleleft$, $\leq$, $\preccurlyeq$, $\ll$, and the other two relations $\parallel$ and $\mathrel{\tikz[baseline=0pt,x=1em,y=1em]{\draw (-.38,.62)--(.38,.25)--(-.43,-.10); \draw (.34,.62)--(-.32,.25)--(.38,-.10);}}$ are symmetric.
% slanted
\margin{No, $\triangleleft$ is not an order relation, it is transitive but not reflexive, for $F\triangleleft F\Longleftrightarrow F\in\unsure{\mathcal{M}|M}$ To have an order relation one would have to write \deleted{$F\triangleleft$} $F\trianglelefteq F\overset{\mathrm{definition}}{\Longleftrightarrow}F\triangleleft F$ \uuline{or} $F=F$.}
The relation $\mathrel{\tikz[baseline=0pt,x=1em,y=1em]{\draw (-.38,.62)--(.38,.25)--(-.43,-.10); \draw (.34,.62)--(-.32,.25)--(.38,-.10);}}$ is \inserted{\unsure{moreover|moreover,}} reflexive, but $\parallel$ is so only if $\unsure{\mathcal{F}|\mathfrak{F}}$ is reduced to its sole \unsure{element|element} $\varnothing_{\mathcal{F}}$, for
\[
 F\parallel F \quad\Longleftrightarrow\quad F=\varnothing_{\mathcal{F}}\qquad \text{empty figure.} \tag*{(1.19)}
\]

\uline{Remarks}\quad a) The relations $\triangleleft$ and $\leq$ determine one another. Knowing $\leq$, one defines $\mathcal{M}\subset\mathcal{F}$ as
\[
\mathcal{M}=\left\{\begin{aligned}
&X\in\mathcal{F}\mid X\;
\underset{\text{for }\leq}{\text{$\mathrm{str}^{\mathrm{t}}$. irreducible}},\;
\text{i.e.}\ X\leq\mathop{\mathrm{Sup}}\limits_i F_i\Rightarrow\\
&\hspace{12em}\exists i\ \text{with }X\leq F_i
\end{aligned}\right\}
\tag*{(1.20)}
\]
then $\triangleleft$ by (1.17).
\deleted{\unsure{And|In}} On the other hand, knowing $\triangleleft$ one defines $\mathcal{M}$ by
\[
\mathcal{M}=\{X\in\mathcal{F}\mid X\triangleleft X\} \tag*{(1.21)}
\]
then, for $F\in\mathcal{F}$, its unfolding
\[
\widetilde F=\{X\in\mathcal{M}\mid X\triangleleft\mathcal{F}\} \tag*{(1.22)}
\]
and finally
\[
F\leq G\quad\Longleftrightarrow\quad\widetilde F\subset\widetilde G \tag*{(1,23)}
\]
Thus, to determine \deleted{the structure, it is \unsure{equivalent|equivalent}} \inserted{the \unsure{minimum|minimal} data which} fix the workshop structure, it is immaterial therein to
